\documentclass[a4paper,oneside,12pt]{amsart}

\usepackage[spanish]{babel}
\usepackage[utf8]{inputenc}
%\usepackage[latin1]{inputenc}
\usepackage{latexsym}
\usepackage{multicol}
\usepackage{enumerate}
\usepackage[heightrounded]{geometry}
\geometry{top=2cm,bottom=2cm,left=2cm,right=2cm}
\usepackage{graphicx}

\DeclareMathOperator{\cond}{cond}
\DeclareMathOperator{\Id}{Id}
\DeclareMathOperator{\re}{Re}
%\DeclareMathOperator{\im}{Im}
\newcommand{\I}{\mathbb{I}}
\newcommand{\R}{\mathbb{R}}
\newcommand{\Q}{\mathbb{Q}}
\newcommand{\C}{\mathbb{C}}
\newcommand{\Z}{\mathbb{Z}}
\newcommand{\N}{\mathbb{N}}
\DeclareMathOperator{\im}{Im}

\newcommand{\nombremateria}
{Investigaci\'on Operativa}
\newcommand{\nombrecuatrimestre}
    {Segundo cuatrimestre de 2020}
\newcommand{\numeroparcial}
    { Recuperatorio del Segundo Parcial }
\newcommand{\fecha}
    {18 de Diciembre de 2020}
\newcommand{\numerotema}
    {\bf 1}

%\oddsidemargin -1cm \evensidemargin -1cm \textwidth 17.5cm
%topmargin 1.2cm \textheight 25cm
%\setlength{\textheight}{25cm}

\def \ds{\displaystyle}
\newtheorem{quest}{Ejercicio}
\thispagestyle{empty}

\begin{document}

\hfill \hfill
\begin{tabular}[c]{|c|c|c|c|c|}
\hline 
1 & 2 & 3 & 4 & Calificaci\'on \\ \hline \hline
\hspace{1.2 cm} & \hspace{1.2 cm} & \hspace{1.2 cm} & \hspace{1.2 cm} &\hspace{1.5 cm} \\
\hspace{1.2 cm} & \hspace{1.2 cm} & \hspace{1.2 cm} & \hspace{1.2 cm} &\hspace{1.5 cm} \\ \hline

\end{tabular}

\vspace{-1cm}

\noindent Nombre y apellido:

\vspace{0.2cm}

\noindent Número de libreta:

\noindent \hrulefill 

\smallskip
\renewcommand{\baselinestretch}{1.1}

\begin{center}
	\large \textbf{\nombremateria} 
	
	\numeroparcial -- \fecha
\end{center}

\vspace{.6cm}
\renewcommand{\theenumi}{\textbf{\arabic{enumi}}}

%%% EMPIEZA EL EJERCICIO 1 %%%

\begin{quest} 
Resolver el siguiente problema de Programación Lineal utilizando el Método de dos Fases de SIMPLEX. Explicitar dónde se realiza el óptimo y el valor que toma la función objetivo en el mismo.
    \[\begin{array}{rrrrrrrr}
        \max & 2x_1 & + & x_2 & - & x_3  \\
        \text{s.a:} & 2x_1 & + & 3x_2 & + & x_3 & \leq & 9 \\
          &  &  & 2x_2 & - & x_3 & \geq & 4 \\
          & x_1  &  &  & - & x_3 & = & 6 \\
         & \multicolumn{7}{l}{x_3\leq 0 \;\;\;x_1,x_2\geq 0}
    \end{array}
    \]

\end{quest}

%%% TERMINA EL EJERCICIO 1 %%%

%%% EMPIEZA EL EJERCICIO 2 %%%

\begin{quest} 
El siguiente diccionario corresponde a una iteración de SIMPLEX de un problema de Programación Lineal con objetivo de maximizar. Sean $a,b,c,d,e,f,g\in \mathbb{R}$, con $b\geq 0$,
decidir y justificar si cada afirmación es verdadera o falsa. Debe considerarse que se utiliza el criterio usual para la selección de la variable que entrará a la base en la siguiente iteración.

\[
\begin{array}{rrrrrrrrr}
    x_3&=&2&-&ax_1&+&x_2&+&dx_5\\
    x_4&=&b&+&cx_1&+&x_2&+&\frac{1}{2}x_5\\
    x_6&=&1&-&x_1&-&x_2&+&fx_5\\ \hline
    z&=&&&ex_1&-&x_2 &+&gx_5
\end{array}
\]

\begin{enumerate}[a)]
    \item Para que la solución básica actual sea degenerada es suficiente pedir que $b=0$.
    \item Para que el problema tenga infinitas soluciones óptimas basta pedir que $b>0$.
    \item Para que la próxima iteración sea degenerada es suficiente que $e>0$ y $b>0$.
    \item Para que en la próxima iteración $x_5$ entre a la base y salga $x_3$ basta pedir que \mbox{$g>0\wedge d<0\wedge f<0$.}
    \item Para que el problema sea no acotado basta pedir que $g>0\;\wedge f\geq 0 \wedge d\geq 0$
\end{enumerate}

\end{quest}

%%% TERMINA EL EJERCICIO 2 %%%

%%% EMPIEZA EL EJERCICIO 3 %%%


\begin{quest} 
Dado el siguiente problema de programación lineal:

\[\begin{array}{rrrrrrrrrrrrrrrr}
    \max & 3x_1 & + & x_2 & - & 2x_3 & + & x_4 & - & x_5 & + & 4x_6 & + & 5x_7 & &\\
    \text{s.a.} & -x_1 & + & 3x_2 & - & x_3 & - & x_4 & + & 2x_5 & + & 6x_6 & - & x_7 & \leq & 5 \\
    &  &  & x_2 & + & 2x_3 & - & 2x_4 & - & x_5 & + & 2x_6 & & &\leq & 5 \\
    & 4x_1 & - & 2x_2 & - & x_3 & + & 3x_4 & + & x_5 & - & x_6 & + &\frac{1}{2}x_7 & \leq & 4 \\
    & 3x_1 & - & x_2  &  &  & + & x_4 &  &  & + & 2x_6 & -&x_7 &\leq & 9 \\
    & \multicolumn{15}{l}{x_1,x_2,x_3,x_4,x_5,x_6,x_7\geq 0}
\end{array}
\]

\begin{enumerate}[a)]
    \item Plantear el problema dual.
    \item Utilizando el problema dual, mostrar que $(\frac{7}{2},5,0,0,0,0,0)$ \textbf{no} es solución óptima del problema primal.
    \item Sabiendo que $x^\ast = (0,17,0,6,0,0,40)$ es solución óptima del problema primal, hallar los valores de $\alpha,\beta, \gamma\in \mathbb{R}$ para los cuales $x^\ast $ es solución óptima del problema con la misma función objetivo, pero cambiando los términos a la derecha de las restricciones por $\alpha$, $\beta$, $\alpha+\beta$, $\alpha+\gamma$, respectivamente.
\end{enumerate}
\end{quest}

%%% TERMINA EL EJERCICIO 3 %%%

%%% EMPIEZA EL EJERCICIO 4 %%%


\begin{quest} 
Utilizar el algoritmo de Branch \& Bound para resolver el siguiente problema de programaci\'on lineal entera, detallando cada divisi\'on en subproblemas durante el algoritmo mediante el esquema de \'arbol visto en clase. Comenzar ramificando en $x_2$ y explicitar el m\'etodo utilizado para la resoluci\'on de cada subproblema (ejemplo: m\'etodo gr\'afico, c\'alculo de v\'ertices de cada regi\'on y posteriormente evaluar en cada uno de ellos, SIMPLEX, etc\'etera). 

\[\begin{array}{rrrrrrrrc}
\max & 3x_1 & + & 4x_2 &  &   \\
\text{s.a:}	 & 2x_1 & + & x_2 & \leq & 14  \\
			 & 2x_1 & + & 3x_2 & \leq & 17 \\
			 &2x_1 & +  & 2x_2 & \leq & 15 \\
			 &	x_1    &  & & \geq & 0 \\
			 & & & x_2 & \geq & 0 \\
			 & \multicolumn{3}{r}{x_1,x_2} & \in & \mathbb{Z} 
\end{array}
\]



\end{quest}

%%% TERMINA EL EJERCICIO 4 %%%

%%% TERMINA EL PARCIAL %%%



\end{document}

